\documentclass[10pt]{article}
\usepackage[T1]{fontenc}
\usepackage[utf8]{inputenc}
\usepackage[french]{babel}
\usepackage{amsmath,amssymb,mathrsfs,mathpazo}
\usepackage{geometry}
\geometry{a4paper,margin=2cm}
\usepackage{graphicx,array,tabularx,tabularray,colortbl,multirow,multicol,booktabs,enumitem,pifont,ulem,fancybox,diagbox,cancel}
\usepackage{pstricks,pstricks-add,pst-plot,pst-tree,pst-node,pst-text,pst-func,pst-math,pst-bspline,pst-3dplot,pst-eucl,pst-solides3d}
\let\pscalculate\undefined
\usepackage{pst-bezier}
\usepackage{tikz,pgfplots,tkz-tab}
\usetikzlibrary{arrows,arrows.meta,patterns,calc,positioning,trees,decorations.markings,shapes.geometric,shapes.misc,shapes.arrows,intersections,decorations.pathmorphing,decorations.pathreplacing,backgrounds,shadows,babel}
\pgfplotsset{compat=1.18}
\usepgfplotslibrary{fillbetween}
\usepackage[np]{numprint}
\usepackage{listings,makecell,pas-tableur,twoopt}
\let\barre\undefined
\usepackage[pstricks]{tablvar}
\definecolor{darkgreen}{RGB}{0,100,0}
\providecommand{\red}{\color{red}}
\providecommand{\blue}{\color{blue}}
\providecommand{\green}{\color{green}}
\providecommand{\black}{\color{black}}
\providecommand{\white}{\color{white}}
\providecommand{\gray}{\color{gray}}
\providecommand{\cyan}{\color{cyan}}
\providecommand{\magenta}{\color{magenta}}
\providecommand{\yellow}{\color{yellow}}
\providecommand{\violet}{\color{violet}}
\providecommand{\orange}{\color{orange}}
\providecommand{\euro}{\texteuro}
\providecommand{\EUR}{\texteuro}
\providecommand{\ieme}{\textsuperscript{e}}
\providecommand{\degres}{\ensuremath{^\circ}}
\providecommand{\degree}{\ensuremath{^\circ}}
\providecommand{\vv}[1]{\overrightarrow{#1}}
\providecommand{\Coord}[1]{#1}
\providecommand{\decosix}{}
\providecommand{\starredbullet}{\ensuremath{\bullet}}
\pagestyle{empty}
\setlength{\parindent}{0pt}
\usepackage{adjustbox}\DeclareUnicodeCharacter{2236}{:}\DeclareUnicodeCharacter{2190}{\ensuremath{\leftarrow}}
\providecommand{\R}{\mathbb{R}}
\providecommand{\N}{\mathbb{N}}
\providecommand{\Z}{\mathbb{Z}}
\providecommand{\Q}{\mathbb{Q}}
\providecommand{\D}{\mathbb{D}}
\providecommand{\C}{\mathbb{C}}
\providecommand{\e}{{\mathrm{e}}}
\providecommand{\vect}[1]{\overrightarrow{#1}}
\providecommand{\vectt}[1]{\overrightarrow{\mathrm{#1}}}
\providecommand{\vvt}[1]{\overrightarrow{\mathrm{#1}}}
\providecommand{\np}[2][]{{#2~#1}}
\providecommand{\vv}[1]{{\overrightarrow{#1}}}
\renewcommand{\d}{{\mathrm{d}}}
\providecommand{\degres}{{^\circ}}
\providecommand{\degre}{{^\circ}}
\providecommand{\blue}{\color{blue}}
\providecommand{\red}{\color{red}}
\providecommand{\enskip}{\quad}
\providecommand{\llbracket}{\lbrack\!\lbrack}
\providecommand{\rrbracket}{\rbrack\!\rbrack}
\providecommand{\euro}{{\text{€}}}
\providecommand{\no}{{\text{n°}}}
\providecommand{\interval}[2]{{[#1;#2]}}
\providecommand{\intervFO}[2]{{[#1;#2[}}
\providecommand{\dotfill}{\ldots}

\providecommand{\ds}{\displaystyle}
\providecommand{\textding}[1]{#1}
\providecommand{\Oij}{\(\mathrm{O};\vec{\imath},\vec{\jmath}\)}
\providecommand{\Oijk}{\(\mathrm{O};\vec{\imath},\vec{\jmath},\vec{k}\)}
\providecommand{\Ouv}{\(\mathrm{O};\vec{u},\vec{v}\)}
\providecommand{\intervFF}[2]{\ensuremath{[#1;#2]}}
\renewcommand{\euro}{\ifmmode\text{\texteuro}\else\texteuro\fi}
\DeclareUnicodeCharacter{00B0}{\ifmmode^{\circ}\else\textdegree\fi}
\DeclareUnicodeCharacter{20AC}{\ifmmode\text{\texteuro}\else\texteuro\fi}
% Énoncés : Éducation nationale. Transcription et corrigés : APMEP, Denis Vergès et contributeurs. Découpage et classement : CoopMaths, Sébastien Lozano, Éric Elter et contributeurs. Adaptation de la présentation : Maths Explorer.
% CC BY-SA : https://coopmaths.fr/www/about/
% https://coopmaths.fr/alea/static/e3c/2020/tex/e3c_2020_00_sujet52_4.tex
\begin{document}
\noindent\textit{APMEP / CoopMaths. CC BY-SA. Présentation adaptée par Maths Explorer.}\par\medskip
\medskip

Le 1\up{er} janvier 2019, le propriétaire d'un appartement a fixé à 650 euros le montant des loyers mensuels pour l'année 2019. Chaque 1\up{er} janvier, le propriétaire augmente de 1,52\,\% le loyer mensuel.

On modélise l'évolution du montant des loyers mensuels par une suite $\left(u_n\right)$. L'arrondi à l'unité du terme $u_n$ représente le montant, en euros, du loyer mensuel fixé le 1\up{er} janvier de l'année $(2019+n)$, pour $n$ entier naturel. Ainsi $u_0 = 650$ euros.

\medskip

\begin{enumerate}
\item 
	\begin{enumerate}
		\item Calculer le montant du loyer mensuel fixé le 1\up{er} janvier 2020.
		\item Quelle est la nature de la suite $\left(u_n\right)$ ? Préciser sa raison et son premier terme.
		\item Calculer le montant du loyer mensuel qui, selon ce modèle, sera fixé pour l'année 2027.
	\end{enumerate}
\item Pour calculer la somme totale des loyers perçus par le propriétaire durant les années 2019
à 2019+A, on utilise la fonction ci-dessous, écrite en langage Python.

\begin{lstlisting}[style=python]
1 def somme(A):
2 	S=0
3 	n=0
4 	while n<=A:
5		S=S+7800*1.0152**n
6 		n = n + 1
7 	return S
\end{lstlisting}

L'exécution de ce programme pour quelques valeurs de A donne les résultats ci-dessous :

\begin{center}

\begin{tabular}[]{|l}
>\! >\! > somme(\textcolor{blue}{0})\\
\textcolor{blue}{7800.0}\\
>\! >\! > somme(\textcolor{blue}{1})\\
\textcolor{blue}{15718.560000000001}\\
>\! >\! > somme(\textcolor{blue}{2})\\
\textcolor{blue}{23757.482112000005}\\
>\! >\! > somme(\textcolor{blue}{3})\\
\textcolor{blue}{31918.595440102407}\\
>\! >\! > somme(\textcolor{blue}{8})\\
\textcolor{blue}{74623.04180934158}\\
\end{tabular}
\end{center}
	\begin{enumerate}
		\item Interpréter, dans le contexte de l'exercice, le résultat obtenu lors de l'appel somme(1).
		\item Déterminer la somme totale des loyers perçus par le propriétaire durant les années
2022 à 2027 incluses. On arrondira le résultat à l'unité.
	\end{enumerate}
\end{enumerate}

\end{document}
